\section{Discussion}
\label{sec:discussion}

This work explored the phenomenon of model collapse, an important concern as AI-generated content permeates the internet and finds its way into future training datasets. Prior work has shown that training on model outputs can lead to degraded performance \citep{martinez2023combining,martinez2023towards,shumailov2023curse,alemohammad2023self,hataya2023will,bertrand2023stability,briesch2023large,dohmatob2024model,dohmatob2024tale}, implying that future model training faces a difficult challenge of ensuring strict training dataset hygiene.
\textbf{For a significantly more thorough discussion of related work, please see Appendix \ref{app:sec:prior_work}.}


Our findings extend these prior works to show that if data \accumulate{accumulates} and models train on a mixture of ``real'' and synthetic data, \accumulate{model collapse no longer occurs}. We show this both experimentally on causal transformers for language modeling, diffusion models for molecule generation, and variational auto-encoders on image data as well as theoretically for linear regression. \textbf{Together, these results strongly suggest that the ``curse of recursion" may not be as dire as had been portrayed \--- provided we accumulate synthetic data alongside real data, rather than replacing real data by synthetic data only.}

Looking to the future, many questions worth investigating remain. For instance, in future work we would like to explore different data generation and accumulation regimes, such as (1) additional ``real'' data being introduced in each model-fitting iteration and (2) different schedules of how much synthetic data is generated at each iteration and (3) human-filtering of generated data, e.g., as done in RLHF. Additionally, we note that in all our experiments, the synthetic dataset is generated by \textit{sampling} from the previous model, i.e., with some stochasticity; in future work, we would like to explore also what happens if data is generated deterministically, e.g. with temperature 0 in a typical language model.

Lastly, it is worth noting that ``model collapse" -- as a term of art -- has been used in various ways by various researchers; so care is required in comparing claims across articles. In reviewing the literature, we identified at least four related phenomena: (0) unbounded test error blowup (as here); (1) modal collapse --- collapse to one (or a few) modes; (2) collapse to uniformity; and (3) amplification of artifacts introduced by models fit to previous synthetic data. Future work should map out what factors cause which to occur and what preventative strategies are effective at addressing each.

\section{Acknowledgements}

The content of this paper does not necessarily reflect the position or the policy of any of the funding agencies/entities. No endorsement should be inferred. M.G. acknowledges support through a grant from the Cooperative AI Foundation. R.S. acknowledges support from Stanford Data Science and from OpenAI's Superalignment Fast Grant Research Fellowship.  A.G. acknowledges support from the NSF CAREER grant DMR-2045181, the Sloan Foundation, and by the Laboratory
for Physical Sciences through the Condensed Matter Theory Center.
D.R. acknowledges support from the National Science Foundation under Cooperative Agreement PHY-2019786 (the NSF \href{http://iaifi.org/}{AI Institute for Artificial Intelligence and Fundamental Interactions}) and appreciates both the sanction and support of Sequoia Capital.
S.K. is partially supported by NSF III 2046795, IIS 1909577, CCF 1934986, NIH 1R01MH116226-01A, NIFA award 2020-67021-32799, the Alfred P. Sloan Foundation, and Google Inc.
